\section{Associated prime cycles and primary decompositions}
\label{section:IV.3}

\oldpage[IV-2]{36}

In this section, we mainly give the translation of results on modules presented in Bourbaki, \emph{Alg. comm.}, chap.~IV, which we follow very closely.
The notions that follow seem to be of interest only in the case of \emph{locally Noetherian} preschemes.

\subsection{Associated prime cycles of a module}
\label{subsection:IV.3.1}

\begin{definition}[3.1.1]
\phantomsection
\label{IV.3.1.1}
Let $X$ be a prescheme and $\sh{F}$ a quasi-coherent $\sh{O}_X$-module.
We say that a point $x\in X$ is \emph{associated} to $\sh{F}$ if the maximal ideal $\mathfrak{m}_x$ of $\sh{O}_x$ is associated to the $\sh{O}_x$-module $\sh{F}_x$ (in other words, if $\mathfrak{m}_x$ is the annihilator of an element of $\sh{F}_x$).
We denote by $\operatorname{Ass}(\sh{F})$ the set of $x\in X$ associated to $\sh{F}$.
We say that a closed irreducible subset $Z$ of $X$ is an \emph{associated prime cycle of $\sh{F}$} if its generic point is associated to $\sh{F}$.
When $\sh{F}=\sh{O}_X$, we also say that the points (resp.\ \emph{prime cycles}) associated to $\sh{F}$ are associated to the \emph{prescheme} $X$.
\end{definition}

We say that an associated prime cycle of $\sh{F}$ (resp.\ $X$) is \emph{embedded} if it is contained in another associated prime cycle of $\sh{F}$ (resp.\ $X$) (in other words, if it is not \emph{maximal} in the set of these cycles).

If $X$ is locally Noetherian, the \emph{irreducible components} of $X$ are none other than the \emph{maximal} (or \emph{non-embedded}) prime cycles associated to $X$.

It is clear that if $x\in\operatorname{Ass}(\sh{F})$, then $\sh{F}_x\neq 0$, in other words
\[
\label{IV.3.1.1.1}
\operatorname{Ass}(\sh{F})\subset\operatorname{Supp}(\sh{F}).
\tag{3.1.1.1}
\]

If $x\in X$ is associated to $\sh{F}$, it is evidently associated to $\sh{F}|U$ for every open neighbourhood $U$ of $x$, and conversely, if it is associated to $\sh{F}|U$ for one of these neighbourhoods, it is associated to $\sh{F}$.

We note finally that for a quasi-coherent $\sh{O}_X$-module $\sh{F}$, the existence of embedded associated prime cycles is a \emph{local} question, for if $y$ and $z$ are two points of $\operatorname{Ass}(\sh{F})$ such that $y\in\overline{\{z\}}$, every neighbourhood of $y$ contains $z$.

\begin{proposition}[3.1.2]
\phantomsection
\label{IV.3.1.2}
Let $A$ be a Noetherian ring, $M$ an $A$-module, $X=\operatorname{Spec}(A)$, $\sh{F}=\widetilde{M}$; for a point $x\in X$ to be associated to $\sh{F}$, it is necessary and sufficient that the prime ideal $\mathfrak{j}_x$ of $A$ be associated to the module $M$ (in other words, be the annihilator of an element $f\in M$).
\end{proposition}

\begin{proof}
This follows from the definition \sref{IV.3.1.1} and from Bourbaki, \emph{loc.\ cit.}, \textsection~1, no~2, cor.\ of prop.~5, applied to $S=A-\mathfrak{j}_x$.
\end{proof}

\begin{proposition}[3.1.3]
\phantomsection
\label{IV.3.1.3}
Let $X$ be a locally Noetherian prescheme, $\sh{F}$ a quasi-coherent $\sh{O}_X$-module, $x$ a point of $X$, and $Z$ the reduced closed sub-prescheme of $X$ having $\overline{\{x\}}$ as underlying space \sref[I]{I.5.2.1}.
The following conditions are equivalent:
\begin{enumerate}
  \item[\emph{a)}] $x\in\operatorname{Ass}(\sh{F})$.
\oldpage[IV-2]{37}
  \item[\emph{b)}] There exist an open neighbourhood $U$ of $x$ and a section $f\in\Gamma(U,\sh{F})$ such that $U\cap Z$ is equal to $\operatorname{Supp}(\sh{O}_U\cdot f)$.
  \item[\emph{b$'$)}] There exist an open neighbourhood $U$ of $x$ and a section $f\in\Gamma(U,\sh{F})$ such that $U\cap Z$ is an irreducible component of $\operatorname{Supp}(\sh{O}_U\cdot f)$.
  \item[\emph{c)}] There exist an open neighbourhood $U$ of $x$ and a submodule of $\sh{F}|U$ isomorphic to $\sh{O}_Z|U$ ($\sh{O}_Z$ being identified with a quotient of $\sh{O}_X$).
  \item[\emph{c$'$)}] There exist an open neighbourhood $U$ of $x$ and a coherent submodule $\sh{G}$ of $\sh{F}|U$ such that $U\cap Z$ is an irreducible component of $\operatorname{Supp}(\sh{G})$.
\end{enumerate}
\end{proposition}

\begin{proof}
It is clear that \emph{c)} implies \emph{b)}, for it suffices to take for $f$ the element of $\Gamma(U,\sh{F})$ which corresponds to the unit section of $\sh{O}_Z|U$.
As $U\cap Z$ is irreducible \sref[0]{0.2.1.6}, \emph{b)} implies \emph{b$'$)}, and \emph{b$'$)} implies \emph{c$'$)} since $\sh{O}_X$ is coherent \sref[0]{0.5.3.4}.
To see that \emph{c$'$)} implies \emph{a)}, we can reduce to the case where $U=X=\operatorname{Spec}(A)$ is affine, $A$ being Noetherian, and where $\sh{F}=\widetilde{M}$, where $M$ is an $A$-module, and $\sh{G}=\widetilde{N}$, where $N$ is a sub-module of $M$, of finite type.
We then know that the minimal elements of $\operatorname{Supp}(\sh{G})$ are the maximal points of $V(\operatorname{Ann}(N))$ \sref[0]{0.1.7.4}, and these are also the minimal elements of $\operatorname{Ass}(N)$ (Bourbaki, \emph{Alg.\ comm.}, chap.~IV, \textsection~1, no~3, cor.~1 of prop.~7);
as $\operatorname{Ass}(N)\subset\operatorname{Ass}(M)=\operatorname{Ass}(\sh{F})$, we see that \emph{c$'$)} implies \emph{a)}.
Finally, \emph{a)} implies \emph{c)} by \sref{IV.3.1.2}, taking $X$ affine again, $\sh{F}=\widetilde{M}$, and $Z$ defined by the ideal $\mathfrak{j}_xA$ (with the notation of \sref{IV.3.1.2}).
\end{proof}

\begin{corollary}[3.1.4]
\phantomsection
\label{IV.3.1.4}
Let $X$ be a locally Noetherian prescheme and $\sh{F}$ a coherent $\sh{O}_X$-module.
Then the maximal prime cycles associated to $\sh{F}$ are the irreducible components of $\operatorname{Supp}(\sh{F})$, and the generic points of these components are the points $x\in X$ such that $\sh{F}_x$ is an $\sh{O}_x$-module of finite nonzero length.
\end{corollary}

\begin{proof}
Indeed, if $x$ is the generic point of an irreducible component $Z$ of $\operatorname{Supp}(\sh{F})$, the equivalence of \emph{a)} and \emph{c$'$)} in \sref{IV.3.1.3} shows that $x$ belongs to $\operatorname{Ass}(\sh{F})$, and $Z$ is an associated prime cycle of $\sh{F}$, necessarily maximal by \sref{IV.3.1.1.1}; the converse follows trivially from \sref{IV.3.1.1.1}. Finally, the last assertion is local and follows from Bourbaki, \emph{Alg.\ comm.}, chap.~IV, \textsection~2, no~5, cor.~2 of prop.~7.
\end{proof}

\begin{corollary}[3.1.5]
\phantomsection
\label{IV.3.1.5}
Let $X$ be a locally Noetherian prescheme and $\sh{F}$ a quasi-coherent $\sh{O}_X$-module.
For $\sh{F}=0$, it is necessary and sufficient that $\operatorname{Ass}(\sh{F})=\varnothing$.
\end{corollary}

\begin{proof}
The question being local, we reduce to the case where $X$ is affine, and the conclusion follows immediately from \sref{IV.3.1.2} and from Bourbaki, \emph{Alg.\ comm.}, chap.~IV, \textsection~1, no~1, cor.~1 of prop.~2.
\end{proof}

\begin{proposition}[3.1.6]
\phantomsection
\label{IV.3.1.6}
Let $X$ be a locally Noetherian prescheme and $\sh{F}$ a coherent $\sh{O}_X$-module; then $\operatorname{Ass}(\sh{F})$ is \emph{locally finite} (that is to say, every point of $X$ admits a neighbourhood whose intersection with $\operatorname{Ass}(\sh{F})$ is finite).
\end{proposition}

\begin{proof}
It suffices to consider the case where $X$ is affine, hence Noetherian, and then the proposition follows from \sref{IV.3.1.2} and from Bourbaki, \emph{Alg.\ comm.}, chap.~IV, \textsection~1, no~4, cor.\ of th.~2.
\end{proof}

\begin{proposition}[3.1.7]
\phantomsection
\label{IV.3.1.7}
Let $X$ be a prescheme.
\begin{enumerate}
  \item[(i)] Let $0\to\sh{F}'\to\sh{F}\to\sh{F}''\to 0$ be an exact sequence of quasi-coherent $\sh{O}_X$-modules.
  Then $\operatorname{Ass}(\sh{F}')\subset\operatorname{Ass}(\sh{F})\subset\operatorname{Ass}(\sh{F}')\cup\operatorname{Ass}(\sh{F}'')$.
  \item[(ii)] Let $\sh{F}$ be a quasi-coherent $\sh{O}_X$-module, $(\sh{F}_\alpha)$ a family of quasi-coherent $\sh{O}_X$-submodules of $\sh{F}$ such that $\sh{F}$ is the union of the $\sh{F}_\alpha$.
  Then $\operatorname{Ass}(\sh{F})=\bigcup_\alpha\operatorname{Ass}(\sh{F}_\alpha)$.
  \item[(iii)] For every family $(\sh{F}_\alpha)$ of quasi-coherent $\sh{O}_X$-modules, we have $\operatorname{Ass}(\bigoplus_\alpha\sh{F}_\alpha)=\bigcup_\alpha\operatorname{Ass}(\sh{F}_\alpha)$.
\end{enumerate}
\end{proposition}

\begin{proof}
We immediately reduce to the corresponding propositions for modules (Bourbaki, \emph{loc.\ cit.}, \textsection~1, no~1, formula~(1), prop.~3 and cor.~1 of prop.~3).
\end{proof}

\oldpage[IV-2]{38}

\begin{proposition}[3.1.8]
\phantomsection
\label{IV.3.1.8}
Let $X$ be a locally Noetherian prescheme, $\sh{F}$ a quasi-coherent $\sh{O}_X$-module, $U$ an open subset of $X$, and $\sh{J}$ a coherent ideal of $\sh{O}_X$ defining a closed sub-prescheme of $X$ having $X-U$ as underlying space.
The following conditions are equivalent:
\begin{enumerate}
  \item[\emph{a)}] $\operatorname{Ass}(\sh{F})\subset U$.
  \item[\emph{b)}] For every affine open subset $V$ of $X$, every section of $\sh{F}$ over $V$ whose restriction to $V\cap U$ is zero is itself zero.
  \item[\emph{c)}] The canonical homomorphism $\sh{F}\to\shHom_{\sh{O}_X}(\sh{J},\sh{F})$ is injective.
\end{enumerate}
\end{proposition}

\begin{proof}
The question being local, we can suppose $X=\operatorname{Spec}(A)$ affine, $A$ being a Noetherian ring, $\sh{F}=\widetilde{M}$, where $M$ is an $A$-module, and $\sh{J}=\widetilde{\mathfrak{J}}$, where $\mathfrak{J}$ is an ideal of $A$.
The homomorphism $M\to\operatorname{Hom}_A(\mathfrak{J},M)$ sends $m\in M$ to the homomorphism $x\mapsto xm$ of $\mathfrak{J}$ into $M$; to say that it is not injective means that there exists $m\neq 0$ in $M$ such that $\mathfrak{J} m=0$.

Let us first show that \emph{c)} implies \emph{a)}. If $\operatorname{Ass}(\sh{F})$ met $X-U$, there would be a prime ideal $\mathfrak{p}\in\operatorname{Ass}(M)$ containing $\mathfrak{J}$, hence an element $m\neq 0$ of $M$ such that $\mathfrak{p}m=0$, and therefore $\mathfrak{J}m=0$, contradicting \emph{c)}.

Next, \emph{b)} implies \emph{c)}. Indeed, if $\mathfrak{J}m=0$ for some $m\neq 0$ in $M$, then for every prime ideal $\mathfrak{q}$ not containing $\mathfrak{J}$ there exists $a\in\mathfrak{J}$ with $a\notin\mathfrak{q}$. The relation $am=0$ then implies that the canonical image of $m$ in $M_\mathfrak{q}$ is zero. Thus $m$ would define a nonzero section of $\sh{F}$ over $X$ whose restriction to $U$ is zero, contrary to \emph{b)}.

Finally, \emph{a)} implies \emph{b)}. The canonical homomorphism
\[
M\longrightarrow\prod_{\mathfrak{p}\in\operatorname{Ass}(M)}M_\mathfrak{p}
\]
is injective: if $N$ is its kernel, then $\operatorname{Ass}(N)\subset\operatorname{Ass}(M)$; if there existed $\mathfrak{p}\in\operatorname{Ass}(N)$, there would be an element $n\in N$ whose annihilator is $\mathfrak{p}$, while the definition of $N$ would give an $s\notin\mathfrak{p}$ such that $sn=0$, a contradiction. Hence $\operatorname{Ass}(N)=\varnothing$, and therefore $N=0$. Under \emph{a)}, a section $m\in M$ whose restriction to $U$ is zero has zero image in $M_\mathfrak{p}$ for every $\mathfrak{p}\in\operatorname{Ass}(M)$; consequently $m=0$.
\end{proof}

\begin{corollary}[3.1.9]
\phantomsection
\label{IV.3.1.9}
Let $X$ be a locally Noetherian prescheme, $\sh{F}$ a quasi-coherent $\sh{O}_X$-module, and $f$ a section of $\sh{O}_X$ over $X$.
For $f$ to be $\sh{F}$-regular \sref[0]{0.15.2.2}, it is necessary and sufficient that $\operatorname{Ass}(\sh{F})\subset X_f$.
\end{corollary}

\begin{proof}
Indeed, it is immediate that to say that $f$ is $\sh{F}$-regular means that the canonical homomorphism $\sh{F}\to\shHom_{\sh{O}_X}(f\sh{O}_X,\sh{F})$ is injective, and it suffices to apply \sref{IV.3.1.8} to the ideal $\sh{J}=f\sh{O}_X$.
\end{proof}

\begin{proposition}[3.1.10]
\phantomsection
\label{IV.3.1.10}
Let $X$ and $Y$ be locally Noetherian preschemes and $f:X\to Y$ an integral morphism.
Then, for every quasi-coherent $\sh{O}_X$-module $\sh{F}$, we have $f(\operatorname{Ass}(\sh{F}))=\operatorname{Ass}(f_*(\sh{F}))$.
\end{proposition}

\begin{proof}
The question being local on $Y$, and the morphism $f$ being affine, we are immediately reduced to the case where $Y$ is affine, in other words, to the following lemma.
\end{proof}

\oldpage[IV-2]{39}

\begin{lemma}[3.1.10.1]
\phantomsection
\label{IV.3.1.10.1}
Let $A$ and $B$ be two Noetherian rings, $\rho:A\to B$ a ring homomorphism making $B$ an integral $A$-algebra, and $M$ a $B$-module.
Then the prime ideals $\mathfrak{p}\in\operatorname{Ass}(M_{[\rho]})$ are the inverse images by $\rho$ of the prime ideals $\mathfrak{q}\in\operatorname{Ass}(M)$.
\end{lemma}

\begin{proof}
Indeed, if $\mathfrak{q}\in\operatorname{Ass}(M)$, then $\mathfrak{q}$ is the annihilator in $B$ of an element $x\in M$, and hence $\rho^{-1}(\mathfrak{q})$ is the annihilator in $A$ of $x$. Conversely, suppose $\mathfrak{p}\in\operatorname{Ass}(M_{[\rho]})$, so that $\mathfrak{p}$ is the inverse image under $\rho$ of the annihilator $\mathfrak{b}$ in $B$ of an element $x\in M$. It follows from the first Cohen--Seidenberg theorem that there exists a prime ideal $\mathfrak{q}$ of $B$ containing $\mathfrak{b}$ and whose inverse image is $\mathfrak{p}$ (Bourbaki, \emph{Alg.\ comm.}, chap.~V, \textsection~2, no~1, cor.~2 of th.~1). By considering a prime ideal minimal among those contained in $\mathfrak{q}$ and containing $\mathfrak{b}$, we may evidently suppose that $\mathfrak{q}$ itself is one of these minimal primes. Since $B.x\subset M$ is isomorphic to $B/\mathfrak{b}$, we then have $\mathfrak{q}\in\operatorname{Ass}(B/\mathfrak{b})\subset\operatorname{Ass}(M)$ (Bourbaki, \emph{Alg.\ comm.}, chap.~IV, \textsection~1, no~4, th.~2).
\end{proof}

\begin{corollary}[3.1.11]
\phantomsection
\label{IV.3.1.11}
Under the hypotheses of \sref{IV.3.1.10}, for $\sh{F}$ to have no embedded associated prime cycle, it suffices that the same be true of $f_*(\sh{F})$.
\end{corollary}

\begin{proof}
Suppose indeed that $f_*(\sh{F})$ has no embedded associated prime cycle.
Note that if $A$ is an integral algebra over a field $k$, all prime ideals of $A$ are maximal (Bourbaki, \emph{Alg.\ comm.}, chap.~V, \textsection~2, no~1, prop.~1); it follows from \sref[I]{I.6.2.2} that the fibres of $f$ are \emph{discrete} spaces. If $x$ and $x'$ are two distinct points of $\operatorname{Ass}(\sh{F})$, neither can be adherent to the other when $f(x)=f(x')$; and if $f(x)\neq f(x')$, then \sref{IV.3.1.10} and the hypothesis imply that neither of the two points $f(x)$ and $f(x')$ can be adherent to the other, so the same is true of $x$ and $x'$.
\end{proof}

\begin{remark}[3.1.12]
\phantomsection
\label{IV.3.1.12}
Under the hypotheses of \sref{IV.3.1.10}, it may happen, on the contrary, that $\sh{F}$ has no embedded associated prime cycle whereas $f_*(\sh{F})$ does.
For example, take $Y=\operatorname{Spec}(k[\mathrm{T}])$, where $k$ is a field (an ``affine line''), and let $X$ be the disjoint union of $X_1=Y$ and $X_2=\operatorname{Spec}(k)$, the morphism $X_2\to Y$ corresponding to the canonical homomorphism $k[\mathrm{T}]\to k[\mathrm{T}]/\mathfrak{m}$, where $\mathfrak{m}$ is the maximal ideal $(\mathrm{T})$.
It is clear that the morphism $f:X\to Y$ is finite. If one takes $\sh{F}=\sh{O}_X$, then $\sh{F}$ has no embedded associated prime cycle, but $f_*(\sh{F})=\widetilde{M}$, where $M$ is the direct sum of the $k[\mathrm{T}]$-modules $k[\mathrm{T}]$ and $k$; hence $\operatorname{Ass}(M)$ consists of the generic point $(0)$ of $Y$ and the point $\mathfrak{m}$.
\end{remark}

\begin{proposition}[3.1.13]
\phantomsection
\label{IV.3.1.13}
Let $X$ be a locally Noetherian prescheme, $U$ an open subset of $X$, and $i:U\to X$ the canonical inclusion.
For every quasi-coherent $\sh{O}_U$-module $\sh{F}$, we have $\operatorname{Ass}(i_*(\sh{F}))=\operatorname{Ass}(\sh{F})$.
\end{proposition}

\begin{proof}
Recall that $i_*(\sh{F})$ is a quasi-coherent $\sh{O}_X$-module \sref[I]{I.2.2} and \sref[I]{I.7.4}; as $i_*(\sh{F})|U=\sh{F}$, we have $\operatorname{Ass}(i_*(\sh{F}))\cap U=\operatorname{Ass}(\sh{F})$, and it therefore remains to prove that $\operatorname{Ass}(i_*(\sh{F}))\subset U$.
But by \sref{IV.3.1.8}, this relation means that for every affine open subset $V$ of $X$, every section of $i_*(\sh{F})$ over $V$ which is zero on $U\cap V$, is zero, a condition which is trivially satisfied since $\Gamma(V,i_*(\sh{F}))=\Gamma(U\cap V,\sh{F})=\Gamma(U\cap V,i_*(\sh{F}))$.
\end{proof}

\oldpage[IV-2]{40}

\subsection{Irredundant decompositions}
\label{subsection:IV.3.2}

\phantomsection
\begin{proposition}[3.2.1]
\label{IV.3.2.1}
Let $X$ be a locally Noetherian prescheme and $U$ a dense open subset of $X$.
The following conditions are equivalent:
\begin{enumerate}
  \item[\emph{a)}] $X$ is reduced.
  \item[\emph{b)}] The induced sub-prescheme on $U$ is reduced, and $X$ has no embedded associated prime cycle.
  \item[\emph{c)}] $X$ has no embedded associated prime cycle and, for every generic point $x$ of an irreducible component of $X$, one has $\operatorname{length}(\sh{O}_x)=1$.
\end{enumerate}
The prime cycles associated with $X$ are then precisely the irreducible components of $X$.
\end{proposition}

\begin{proof}
It is clear that if $X$ is reduced, then so is the sub-prescheme induced on $U$.
Moreover, since the existence of embedded associated prime cycles is a local question, we may restrict to the case where $X=\operatorname{Spec}(A)$ is affine and $A$ is Noetherian.
If $A$ is reduced, its minimal prime ideals form a reduced primary decomposition of $(0)$ (Bourbaki, \emph{Alg.\ comm.}, chap.~IV, \textsection~2, no~5, prop.~10) and are precisely the elements of $\operatorname{Ass}(A)$; hence $A$ has no embedded associated prime ideals.  Thus \emph{a)} implies \emph{b)}.
It is immediate that \emph{b)} implies \emph{c)}, since a generic point $x$ of an irreducible component belongs to $U$, and consequently $\sh{O}_x$ is a field.
Finally, \emph{c)} implies \emph{a)}.  Indeed, let $\sh{N}$ be the nilradical of $\sh{O}_X$, which is a coherent ideal.  By hypothesis, $\operatorname{Supp}(\sh{N})$ contains none of the generic points of the irreducible components of $X$.  If $\operatorname{Supp}(\sh{N})$ were nonempty and $x$ were one of its maximal points, criterion \sref{IV.3.1.3}, \emph{c$'$)}, would give $x\in\operatorname{Ass}(\sh{O}_X)$; then $\overline{\{x\}}$ would be an embedded associated prime cycle of $X$, contrary to the hypothesis.  Hence $\sh{N}=0$.
\end{proof}

\phantomsection
\begin{definition}[3.2.2]
\label{IV.3.2.2}
Let $X$ be a locally Noetherian prescheme and $\sh{F}$ a coherent $\sh{O}_X$-module.
We say that $\sh{F}$ is \emph{reduced} if it satisfies the following two conditions: 1\textsuperscript{o} $\sh{F}$ has no embedded associated prime cycle; 2\textsuperscript{o} for every maximal point $x$ of $\operatorname{Supp}(\sh{F})$, one has $\operatorname{length}(\sh{F}_x)=1$.
\end{definition}

Condition 1\textsuperscript{o} means that the prime cycles associated with $\sh{F}$ are the irreducible components of $\operatorname{Supp}(\sh{F})$ \sref{IV.3.1.4}, and condition 2\textsuperscript{o} means that for every generic point $x$ of such a component one has $\operatorname{length}(\sh{F}_x)=1$.

For an affine scheme $X$, this definition gives in particular the notion of a \emph{reduced module} over a Noetherian ring $A$: a finitely generated $A$-module $M$ is said to be \emph{reduced} if it has no embedded associated prime ideals and if, for every $\mathfrak{p}\in\operatorname{Ass}(M)$, one has $\operatorname{length}_{A_\mathfrak{p}}(M_\mathfrak{p})=1$.
Returning to a locally Noetherian prescheme $X$ and a coherent $\sh{O}_X$-module $\sh{F}$, we say that $\sh{F}$ is \emph{reduced at a point} $x\in X$ if $\sh{F}_x$ is a reduced $\sh{O}_x$-module.  Equivalently, on the local scheme $\operatorname{Spec}(\sh{O}_x)$, the module associated with $\sh{F}_x$ is reduced.  Thus $x$ belongs to no embedded associated prime cycle of $\sh{F}$ and $\operatorname{length}(\sh{F}_z)=1$ \emph{for every maximal point $z$ of $\operatorname{Supp}(\sh{F})$ such that $x\in\overline{\{z\}}$}.
It is clear that if $\sh{F}$ is a reduced coherent $\sh{O}_X$-module, then it is reduced at every point of $X$.  Conversely, if $\sh{F}$ is \emph{reduced at a point} $x$, there exists an open neighbourhood $U$ of $x$ such that $\sh{F}|U$ is a reduced $\sh{O}_U$-module: it suffices to take $U$ meeting no embedded associated prime cycle of $\sh{F}$ (such a neighbourhood exists because these cycles form a locally finite family of closed subsets of $X$).
To say that $\sh{O}_X$ is reduced at $x$ is equivalent to saying that $X$ is \emph{reduced at the point} $x$.

\oldpage[IV-2]{41}

\phantomsection
\begin{proposition}[3.2.3]
\label{IV.3.2.3}
Let $X$ be a locally Noetherian prescheme, $U$ an open subset of $X$, and $\sh{F}$ a coherent $\sh{O}_X$-module such that $U\cap\operatorname{Supp}(\sh{F})$ is dense in $\operatorname{Supp}(\sh{F})$.
The following conditions are equivalent:
\begin{enumerate}
  \item[\emph{a)}] $\sh{F}$ is reduced.
  \item[\emph{b)}] $\sh{F}|U$ is reduced, and $\sh{F}$ has no embedded associated prime cycle.
  \item[\emph{c)}] There exist a reduced closed sub-prescheme $X'$ of $X$ and a coherent torsion-free $\sh{O}_{X'}$-module $\sh{F}'$ of rank $1$ on every irreducible component of $X'$ such that, if $j:X'\to X$ is the canonical inclusion, then $j_*(\sh{F}')=\sh{F}$.
\end{enumerate}
Moreover, when these conditions hold, $X'$ is defined by the ideal $\sh{J}$ of $\sh{O}_X$ that annihilates $\sh{F}$.
\end{proposition}

\begin{proof}
To see that \emph{c)} implies \emph{a)}, we may, by \sref{IV.3.1.3}, restrict to the case where $X'=X$ is integral, with generic point $x$.  We may further suppose that $X=\operatorname{Spec}(A)$ is affine and $\sh{F}=\widetilde{M}$, where $A$ is integral and $M$ is a torsion-free $A$-module of rank $1$.  The annihilator of every element of $M$ is then zero, so $\operatorname{Ass}(\sh{F})=\{x\}$ and $\sh{F}_x$ is isomorphic to $\sh{O}_x$, the field of fractions of $A$; hence the conditions of Definition \sref{IV.3.2.2} are satisfied.
Since the existence of embedded associated prime cycles is a local question, if $\sh{F}$ has no such cycles and $U\cap\operatorname{Supp}(\sh{F})$ is dense in $\operatorname{Supp}(\sh{F})$, then $\operatorname{Ass}(\sh{F})=\operatorname{Ass}(\sh{F}|U)$.  Thus \emph{a)} and \emph{b)} are equivalent.
If \emph{a)} holds, take for $X'$ the reduced closed sub-prescheme of $X$ whose underlying space is $\operatorname{Supp}(\sh{F})$ \sref[I]{I.5.2.1}, and put $\sh{F}'=j^*(\sh{F})$.  A point $x$ of $\operatorname{Ass}(\sh{F})$ is necessarily a maximal point of $X'$, and since $\operatorname{length}(\sh{F}_x)=1$, $\sh{F}'_x$ is isomorphic to $\sh{F}_x$, hence to the residue field $\boldsymbol{k}(x)=\sh{O}_{X',x}$.  This proves that $\sh{F}'$ is torsion-free and of rank $1$ \sref[I]{I.7.4.6} and \sref[I]{I.7.4.2}.
Finally, the last assertion is immediate, since for every $y\in X'$ the annihilator of the $\sh{O}_{X',y}$-module $\sh{F}'_y$ is zero.
\end{proof}

\phantomsection
\begin{definition}[3.2.4]
\label{IV.3.2.4}
Let $X$ be a locally Noetherian prescheme and $\sh{F}$ a quasi-coherent $\sh{O}_X$-module.
We say that $\sh{F}$ is \emph{irredundant} if $\operatorname{Ass}(\sh{F})$ consists of a single element $x$; if $\sh{F}$ is of finite type, we say that $\sh{F}$ is \emph{integral} if, in addition, $\sh{F}$ is reduced (in other words, if $\operatorname{length}(\sh{F}_x)=1$).
We say that a quasi-coherent $\sh{O}_X$-submodule $\sh{G}$ of a quasi-coherent $\sh{O}_X$-module $\sh{F}$ is \emph{primary in $\sh{F}$} if $\sh{F}/\sh{G}$ is irredundant.
\end{definition}

For an affine scheme $X$, this definition gives in particular the notion of an \emph{integral module} over a Noetherian ring $A$: an $A$-module $M$ is said to be \emph{integral} if it is of finite type, if $M$ is primary (that is, if $\operatorname{Ass}(M)$ consists of a single prime ideal $\mathfrak{p}$), and if, in addition, $\operatorname{length}_{A_\mathfrak{p}}(M_\mathfrak{p})=1$.
Returning to an arbitrary locally Noetherian prescheme $X$ and a coherent $\sh{O}_X$-module $\sh{F}$, we say that $\sh{F}$ is \emph{integral at a point} $x\in X$ if $\sh{F}_x$ is an integral $\sh{O}_x$-module.  Equivalently, $x$ belongs to only one associated prime cycle of $\sh{F}$ (necessarily a nonembedded one), and $\operatorname{length}(\sh{F}_z)=1$ at its generic point $z$.
It is clear that if $\sh{F}$ is an integral coherent $\sh{O}_X$-module, then it is integral at every point of $X$.  Conversely, if $\sh{F}$ is \emph{integral at a point} $x$, there exists an open neighbourhood $U$ of $x$ such that $\sh{F}|U$ is an integral $\sh{O}_U$-module: it suffices to take $U$ such that $\sh{F}|U$ is a reduced $\sh{O}_U$-module \sref{IV.3.2.2}.

\oldpage[IV-2]{42}

We say that a locally Noetherian prescheme $X$ is \emph{irredundant} if $\sh{O}_X$ is irredundant (which implies that $X$ is \emph{irreducible}).  For $X$ to be \emph{integral}, it is necessary and sufficient that the $\sh{O}_X$-module $\sh{O}_X$ be integral \sref[I]{I.2.1.8} and \sref{IV.3.2.1}.
If $\sh{O}_X$ is integral at a point $x$, that is, if the ring $\sh{O}_x$ is integral, we say that $X$ is \emph{integral at the point} $x$.
We say that a closed sub-prescheme $Y$ of $X$ is \emph{primary in $X$} if the ideal $\sh{J}$ of $\sh{O}_X$ defining $Y$ is primary in $\sh{O}_X$.

\phantomsection
\begin{definition}[3.2.5]
\label{IV.3.2.5}
Let $X$ be a locally Noetherian prescheme and $\sh{F}$ a coherent $\sh{O}_X$-module.
An \emph{irredundant decomposition} of $\sh{F}$ is a family $(\sh{F}_\alpha)_{\alpha\in I}$ of quotient $\sh{O}_X$-modules of $\sh{F}$ such that the $\sh{F}_\alpha$ are irredundant, the family $(\operatorname{Supp}(\sh{F}_\alpha))_{\alpha\in I}$ is locally finite, and the canonical homomorphism $\sh{F}\to\bigoplus_{\alpha\in I}\sh{F}_\alpha$ is injective.
We say that such a decomposition is \emph{reduced} if the sets $\operatorname{Ass}(\sh{F}_\alpha)$ are pairwise distinct and there is no proper subset $J\subsetneq I$ for which the subfamily $(\sh{F}_\alpha)_{\alpha\in J}$ is also an irredundant decomposition of $\sh{F}$.
\end{definition}

If $(\sh{F}_\alpha)_{\alpha\in I}$ is an irredundant decomposition (resp.\ a reduced irredundant decomposition) of $\sh{F}$ and $\sh{F}_\alpha=\sh{F}/\sh{G}_\alpha$, we also say that the family $(\sh{G}_\alpha)_{\alpha\in I}$ of $\sh{O}_X$-submodules of $\sh{F}$ is a \emph{primary decomposition of $0$ in $\sh{F}$}.  The injectivity of the canonical homomorphism $\sh{F}\to\bigoplus_{\alpha\in I}\sh{F}_\alpha$ is equivalent to the condition $\bigcap_{\alpha\in I}\sh{G}_\alpha=0$.

If $(\sh{F}_\alpha)$ is an irredundant decomposition of $\sh{F}$, to say that it is reduced is equivalent to saying that the sets $\operatorname{Ass}(\sh{F}_\alpha)$ are pairwise distinct and contained in $\operatorname{Ass}(\sh{F})$.  If $\operatorname{Ass}(\sh{F}_\alpha)=\{x_\alpha\}$ for every $\alpha\in I$, then $\alpha\mapsto x_\alpha$ is a bijection from $I$ onto $\operatorname{Ass}(\sh{F})$.  These properties are local and therefore follow from Bourbaki, \emph{Alg.\ comm.}, chap.~IV, \textsection~2, no~3, prop.~4.

\phantomsection
\begin{proposition}[3.2.6]
\label{IV.3.2.6}
Let $X$ be a locally Noetherian prescheme and $\sh{F}$ a coherent $\sh{O}_X$-module.
Then there exists a reduced irredundant decomposition $(\sh{F}^{(x)})_{x\in\operatorname{Ass}(\sh{F})}$ formed of coherent $\sh{O}_X$-modules such that, for every $x\in\operatorname{Ass}(\sh{F})$, one has $\operatorname{Ass}(\sh{F}^{(x)})=\{x\}$.
For every $x\in\operatorname{Ass}(\sh{F})$ such that $\overline{\{x\}}$ is not embedded, $\sh{F}^{(x)}$ is uniquely determined as the image of the canonical homomorphism $\sh{F}\to i_*(i^*(\sh{F}))$, where $i$ is the canonical morphism $\operatorname{Spec}(\boldsymbol{k}(x))\to X$.
\end{proposition}

\begin{proof}
For every $x\in\operatorname{Ass}(\sh{F})$, let $U$ be an affine open neighbourhood of $x$, with ring $A$, and write $\sh{F}|U=\widetilde{M}$, where $M$ is a finitely generated $A$-module.
We know (Bourbaki, \emph{Alg.\ comm.}, chap.~IV, \textsection~1, no~1, prop.~4) that there exists a submodule $N$ of $M$ such that, on setting $P=M/N$, one has $\operatorname{Ass}(P)=\{x\}$ and $\operatorname{Ass}(N)=\operatorname{Ass}(M)-\{x\}$.
Let $\sh{G}=\widetilde{P}$, a quasi-coherent $\sh{O}_U$-module, and let $j:U\to X$ be the canonical inclusion.  Let $u:j^*(\sh{F})\to\sh{G}$ be the surjective homomorphism corresponding to $M\to P$.  It induces a homomorphism $j_*(u):j_*(j^*(\sh{F}))\to j_*(\sh{G})$, and hence, by composition, a homomorphism
\[
v:\sh{F}\xrightarrow{\rho_\sh{F}}j_*(j^*(\sh{F}))\xrightarrow{j_*(u)}j_*(\sh{G})
\]
whose restriction to $U$ is $u$.  We denote by $\sh{F}^{(x)}$ the image of $\sh{F}$ under this homomorphism; it is a coherent $\sh{O}_X$-module \sref[I]{I.6.1.1}.
By \sref{IV.3.1.13}, one has $\operatorname{Ass}(j_*(\sh{G}))=\{x\}$, and therefore, by \sref{IV.3.1.7}, $\operatorname{Ass}(\sh{F}^{(x)})=\{x\}$, since $\sh{F}^{(x)}\neq 0$.

Let $\sh{N}^{(x)}=\operatorname{Ker}(v)$.  Then $\sh{N}^{(x)}|U=\widetilde{N}$, and hence $x\notin\operatorname{Ass}(\sh{N}^{(x)})$.
It follows that the homomorphism $\sh{F}\to\bigoplus_{x\in\operatorname{Ass}(\sh{F})}\sh{F}^{(x)}$ is injective.  Indeed, its kernel $\sh{H}$ is contained in every $\sh{N}^{(x)}$, so $\operatorname{Ass}(\sh{H})$ is contained in the intersection of the sets $\operatorname{Ass}(\sh{N}^{(x)})$, which is empty; consequently, by \sref{IV.3.1.5}, $\sh{H}=0$.
Since $\operatorname{Ass}(\sh{F})$ is locally finite \sref{IV.3.1.6}, it is clear that $(\sh{F}^{(x)})_{x\in\operatorname{Ass}(\sh{F})}$ is a reduced irredundant decomposition of $\sh{F}$ satisfying the conditions of the statement.
Finally, when $\overline{\{x\}}$ is not embedded, the characterization of $\sh{F}^{(x)}$ follows from Bourbaki, \emph{Alg.\ comm.}, chap.~IV, \textsection~2, no~3, prop.~5, since the question is local, together with \sref[I]{I.1.6.7}.
\end{proof}

\oldpage[IV-2]{43}

\phantomsection
\begin{corollary}[3.2.7]
\label{IV.3.2.7}
Under the hypotheses of \sref{IV.3.2.6}, if $\sh{F}$ has no embedded associated prime cycle, there exists a unique reduced irredundant decomposition of $\sh{F}$.
\end{corollary}

\phantomsection
\begin{corollary}[3.2.8]
\label{IV.3.2.8}
Let $X$ be a Noetherian prescheme and $\sh{F}$ a coherent $\sh{O}_X$-module.
There exists a finite filtration $(\sh{F}_i)_{0\leqslant i\leqslant n}$ of $\sh{F}$ such that $\sh{F}_0=\sh{F}$ and $\sh{F}_n=0$, formed of coherent $\sh{O}_X$-modules, and such that the quotients $\sh{F}_i/\sh{F}_{i+1}$ are either zero or irredundant, with $\operatorname{Ass}(\sh{F}_i/\sh{F}_{i+1})\subset\operatorname{Ass}(\sh{F})$.
\end{corollary}

\begin{proof}
Indeed, by \sref{IV.3.2.6}, $\sh{F}$ is isomorphic to an $\sh{O}_X$-submodule of a finite direct sum $\bigoplus_{j=1}^n\sh{G}_j$, where the $\sh{G}_j$ are irredundant and coherent.  Since every quasi-coherent $\sh{O}_X$-submodule of $\sh{G}_j$ is either zero or irredundant \sref{IV.3.1.7}, the modules
\[
\sh{F}_i=\sh{F}\cap\left(\bigoplus_{j=1}^{n-i}\sh{G}_j\right)
\]
satisfy the requirements, since $\sh{F}_i/\sh{F}_{i+1}$ is isomorphic to a coherent $\sh{O}_X$-submodule of $\sh{G}_{n-i}$.
\end{proof}

\subsection{Relations with flatness}
\label{subsection:IV.3.3}

\phantomsection
\begin{proposition}[3.3.1]
\label{IV.3.3.1}
Let $f:X\to Y$ be a morphism, $\sh{F}$ a quasi-coherent $\sh{O}_X$-module that is $f$-flat, and $\sh{E}$ a quasi-coherent $\sh{O}_Y$-module.
If, for every $y\in Y$, we set $\sh{F}_y=\sh{F}\otimes_{\sh{O}_Y}\boldsymbol{k}(y)$, then
\[
\label{IV.3.3.1.1}
\operatorname{Ass}(\sh{E}\otimes_{\sh{O}_Y}\sh{F})\supset\bigcup_{y\in\operatorname{Ass}(\sh{E})}\operatorname{Ass}(\sh{F}_y)
\tag{3.3.1.1}
\]
and the two sides are equal if $Y$ is locally Noetherian.
\end{proposition}

\begin{proof}
(Of course, $\sh{F}_y$ is a sheaf on the fibre $f^{-1}(y)$, and this fibre is identified with a subspace of $X$ \sref[I]{I.3.6.1}.)
The question being local on $X$ and on $Y$, we reduce to the case where $X$ and $Y$ are affine, and the proposition is then proved in Bourbaki, \emph{Alg.\ comm.}, chap.~IV, \textsection~2, no~6, th.~2.
\end{proof}

\phantomsection
\begin{corollary}[3.3.2]
\label{IV.3.3.2}
Let $Y$ be a locally Noetherian prescheme with no embedded associated prime cycles, $f:X\to Y$ a morphism, and $\sh{F}$ a quasi-coherent $\sh{O}_X$-module that is $f$-flat.
Then, for every $x\in\operatorname{Ass}(\sh{F})$, $f(x)$ is a maximal point of $Y$.
\end{corollary}

\begin{proof}
It suffices to apply \sref{IV.3.3.1} with $\sh{E}=\sh{O}_Y$, since $\operatorname{Ass}(\sh{O}_Y)$ is by hypothesis the set of maximal points of $Y$.
\end{proof}

\phantomsection
\begin{corollary}[3.3.3]
\label{IV.3.3.3}
Under the hypotheses of \sref{IV.3.3.1}, suppose in addition that $X$ and $Y$ are locally Noetherian and that $\sh{E}$ and $\sh{F}$ are coherent.
Then the following conditions are equivalent:
\begin{enumerate}
  \item[\emph{a)}] $\sh{E}\otimes_Y\sh{F}$ has no embedded associated prime cycle.
  \item[\emph{b)}] For every point $y\in\operatorname{Ass}(\sh{E})\cap f(\operatorname{Supp}(\sh{F}))$, $\overline{\{y\}}$ is a nonembedded associated prime cycle of $\sh{E}$ and $\sh{E}_y\otimes_{\boldsymbol{k}(y)}\sh{F}_y$ has no embedded associated prime cycle.
\end{enumerate}
\end{corollary}

\begin{proof}
\oldpage[IV-2]{44}
Suppose \emph{a)} is satisfied.
The hypotheses imply that $\sh{E}\otimes_Y\sh{F}$ is a coherent $\sh{O}_X$-module \sref[0]{0.5.3.11} and \sref[0]{0.5.3.5}; its associated prime cycles are therefore the irreducible components of $\operatorname{Supp}(\sh{E}\otimes_Y\sh{F})$ \sref{IV.3.1.4}, and for every maximal point $x$ of $\operatorname{Supp}(\sh{E}\otimes_Y\sh{F})$, $f(x)=y$ is a maximal point of $\operatorname{Supp}(\sh{E})$ and $x$ is a maximal point of $\operatorname{Supp}(\sh{F}_y)$ \sref{IV.2.5.5}.
As, by virtue of \sref{IV.3.3.1} and the fact that $y\in f(\operatorname{Supp}(\sh{F}))$ implies $\sh{F}_y\neq 0$ \sref[I]{I.9.1.13}, every point of $\operatorname{Ass}(\sh{E})\cap f(\operatorname{Supp}(\sh{F}))$ is the image by $f$ of a maximal point of $\operatorname{Supp}(\sh{E}\otimes_Y\sh{F})$, we see that condition \emph{b)} is satisfied.

Conversely, suppose \emph{b)} is satisfied, and let us show that if $z$, $z'$ are two distinct points of $\operatorname{Ass}(\sh{E}\otimes_Y\sh{F})$, neither of the two can be adherent to the other.
In the first place, if $f(z)=f(z')=y$, we have $y\in\operatorname{Ass}(\sh{E})$ and $z$ and $z'$ belong to $\operatorname{Ass}(\sh{F}_y)$ by \sref{IV.3.3.1}, whence $y\in f(\operatorname{Supp}(\sh{F}))$; as by hypothesis neither of the two points $z$, $z'$ is adherent to the other \emph{within} $f^{-1}(y)$, neither can be adherent to the other \emph{within} $X$.
If $y=f(z)$ and $y'=f(z')$ are distinct, they belong to $\operatorname{Ass}(\sh{E})\cap f(\operatorname{Supp}(\sh{F}))$, hence neither of the two can be adherent to the other in $Y$; it follows from the continuity of $f$ that neither of the points $z$, $z'$ can be adherent to the other in $X$.
\end{proof}

\phantomsection
\begin{proposition}[3.3.4]
\label{IV.3.3.4}
Let $X$ and $Y$ be locally Noetherian preschemes, $f:X\to Y$ a morphism, $\sh{E}$ a coherent $\sh{O}_Y$-module, and $\sh{F}$ a coherent $\sh{O}_X$-module that is $f$-flat.
Then the following conditions are equivalent:
\begin{enumerate}
  \item[\emph{a)}] $\sh{E}\otimes_Y\sh{F}$ is reduced \sref{IV.3.2.2}.
  \item[\emph{b)}] For every point $y\in\operatorname{Ass}(\sh{E})\cap f(\operatorname{Supp}(\sh{F}))$, $\overline{\{y\}}$ is a nonembedded associated prime cycle of $\sh{E}$, $\operatorname{length}(\sh{E}_y)=1$, and $\sh{F}_y$ is reduced.
\end{enumerate}
\end{proposition}

\begin{proof}
Suppose \emph{a)} is satisfied.
We already know from \sref{IV.3.3.3} that, for every $y\in\operatorname{Ass}(\sh{E})$ lying in $f(\operatorname{Supp}(\sh{F}))$, $\overline{\{y\}}$ is a nonembedded associated prime cycle of $\sh{E}$ and $\sh{F}_y$ has no embedded associated prime cycle.
Furthermore, by \sref{IV.2.5.5}, for every $x\in\operatorname{Ass}(\sh{E}\otimes_Y\sh{F})\cap f^{-1}(y)$, we have
\[
1=\operatorname{length}((\sh{E}\otimes_Y\sh{F})_x)
=\operatorname{length}(\sh{E}_y)\operatorname{length}((\sh{F}_y)_x),
\]
hence $\operatorname{length}(\sh{E}_y)=\operatorname{length}((\sh{F}_y)_x)=1$, which proves \emph{b)}.

Conversely, suppose \emph{b)} is satisfied; we already know that every point $x\in\operatorname{Ass}(\sh{E}\otimes_Y\sh{F})$ is a maximal point of $\operatorname{Supp}(\sh{E}\otimes_Y\sh{F})$, that $y=f(x)$ is a maximal point of $\operatorname{Supp}(\sh{E})$ and $x$ a maximal point of $\operatorname{Supp}(\sh{F}_y)$ \sref{IV.3.3.1} and \sref{IV.3.3.3}; in addition it follows from the hypothesis and from \sref{IV.2.5.5} that $\operatorname{length}((\sh{E}\otimes_Y\sh{F})_x)=1$, which proves \emph{a)}.
\end{proof}

\phantomsection
\begin{corollary}[3.3.5]
\label{IV.3.3.5}
Let $X$ and $Y$ be locally Noetherian preschemes and let $f:X\to Y$ be a flat morphism.
If $Y$ is reduced at the points of $f(X)$ and $f^{-1}(y)$ is a reduced $\boldsymbol{k}(y)$-prescheme for every $y\in f(X)$, then $X$ is reduced.
\end{corollary}

\begin{proof}
Since the nilradical $\sh{N}_Y$ is coherent, the set of points where $Y$ is reduced is open \sref[0]{0.5.2.2}, and we may restrict to the case where $Y$ is reduced.
It suffices then to apply \sref{IV.3.3.4} with $\sh{E}=\sh{O}_Y$ and $\sh{F}=\sh{O}_X$.
\end{proof}

\phantomsection
\begin{proposition}[3.3.6]
\label{IV.3.3.6}
Let $f:X\to S$ and $g:Y\to S$ be two morphisms, $\sh{F}$ a quasi-coherent $\sh{O}_X$-module, and $\sh{G}$ a quasi-coherent $\sh{O}_Y$-module.
Suppose that: 1\textsuperscript{o} $\sh{G}$ is $g$-flat; 2\textsuperscript{o} $X$ is locally Noetherian, and for every $s\in S$, $g^{-1}(s)$ is locally Noetherian (which will hold if $Y$ is also locally Noetherian).
Let $Z=X\times_S Y$; for every pair $(x,y)$ such that $x\in X$, $y\in Y$,
\oldpage[IV-2]{45}
and $f(x)=g(y)=s$, let $\mathrm{T}_{x,y}$ be the prescheme $\operatorname{Spec}(\boldsymbol{k}(x)\otimes_{\boldsymbol{k}(s)}\boldsymbol{k}(y))$, and let $\mathrm{I}_{x,y}$ be the image of $\operatorname{Ass}(\sh{O}_{\mathrm{T}_{x,y}})$ by the canonical monomorphism $\mathrm{T}_{x,y}\to Z$ \sref[I]{I.3.4.9}.
Then
\[
\label{IV.3.3.6.1}
\operatorname{Ass}(\sh{F}\otimes_S\sh{G})=\bigcup_{x\in\operatorname{Ass}(\sh{F})}\bigg(\bigcup_{y\in\operatorname{Ass}(\sh{G}_{f(x)})}\mathrm{I}_{x,y}\bigg)
\tag{3.3.6.1}
\]
where for all $s\in S$, $\sh{G}_s=\sh{G}\otimes_{\sh{O}_S}\boldsymbol{k}(s)$.
\end{proposition}

\begin{proof}
Let $p:Z\to X$, $q:Z\to Y$ be the canonical projections, so that we have the commutative diagram
\[
\xymatrix{
X\ar[d]_{f} & Z\ar[l]_{p}\ar[d]^{q}\\
S & Y\ar[l]^{g}
}
\]
Set $\sh{G}'=q^*(\sh{G})$, so that $\sh{F}\otimes_S\sh{G}=\sh{F}\otimes_X\sh{G}'$; as $\sh{G}'$ is $p$-flat \sref{IV.2.1.4}, it follows from \sref{IV.3.3.1} that
\[
\label{IV.3.3.6.2}
\operatorname{Ass}(\sh{F}\otimes_S\sh{G})=\bigcup_{x\in\operatorname{Ass}(\sh{F})}\operatorname{Ass}(\sh{G}'_x)
\tag{3.3.6.2}
\]
with $\sh{G}'_x=\sh{G}'\otimes_{\sh{O}_X}\boldsymbol{k}(x)$.
If $s=f(x)$, we have $\sh{G}'_x=\sh{G}_s\otimes_{\boldsymbol{k}(s)}\boldsymbol{k}(x)$, and $p^{-1}(x)=g^{-1}(s)\otimes_{\boldsymbol{k}(s)}\boldsymbol{k}(x)$; furthermore, as the field $\boldsymbol{k}(x)$ is a flat $\boldsymbol{k}(s)$-module, the morphism $p^{-1}(x)\to g^{-1}(s)$ is flat; applying \sref{IV.3.3.1} to this morphism, we get
\[
\label{IV.3.3.6.3}
\operatorname{Ass}(\sh{G}'_x)=\bigcup_{y\in\operatorname{Ass}(\sh{G}_s)}\operatorname{Ass}(\sh{O}_{\mathrm{T}_{x,y}})
\tag{3.3.6.3}
\]
whence the proposition.
\end{proof}

We note that if, in the statement, we suppress hypothesis 2\textsuperscript{o}, we can still conclude, by virtue of \sref{IV.3.3.1}, the relation
\[
\label{IV.3.3.6.4}
\operatorname{Ass}(\sh{F}\otimes_S\sh{G})\supset\bigcup_{x\in\operatorname{Ass}(\sh{F})}\bigg(\bigcup_{y\in\operatorname{Ass}(\sh{G}_{f(x)})}\mathrm{I}_{x,y}\bigg).
\tag{3.3.6.4}
\]

\phantomsection
\begin{corollary}[3.3.7]
\label{IV.3.3.7}
Under the hypotheses of \sref{IV.3.3.6}, suppose in addition that $S$ is also locally Noetherian and that $f(\operatorname{Ass}(\sh{F}))\subset\operatorname{Ass}(\sh{O}_S)$.
Then
\[
\label{IV.3.3.7.1}
\operatorname{Ass}(\sh{F}\otimes_S\sh{G})=\bigcup_{(x,y)\in C}\mathrm{I}_{x,y}
\tag{3.3.7.1}
\]
where $C$ is the set of pairs $(x,y)$ such that $x\in\operatorname{Ass}(\sh{F})$, $y\in\operatorname{Ass}(\sh{G})$ and $f(x)=g(y)$.
\end{corollary}

\begin{proof}
Since $\sh{G}$ is $g$-flat, it follows from \sref{IV.3.3.1} that the relation
``$s\in\operatorname{Ass}(\sh{O}_S)$ and $y\in\operatorname{Ass}(\sh{G}_s)$''
is equivalent to ``$y\in\operatorname{Ass}(\sh{G})$''; the conclusion follows from \sref{IV.3.3.6.1}.
\end{proof}

\phantomsection
\begin{remarks}[3.3.8]
\label{IV.3.3.8}
We shall see later \sref{IV.4.2.2} that under the hypotheses of \sref{IV.3.3.6}, $\mathrm{T}_{x,y}$ is a prescheme \emph{with no embedded associated prime cycle}; it will follow that if $\sh{F}$ and the $\sh{G}_s$ have no embedded associated prime cycles, the same is true of $\sh{F}\otimes_S\sh{G}$.
\end{remarks}

\phantomsection
\begin{corollary}[3.3.9]
\label{IV.3.3.9}
Under the conditions of \sref{IV.3.3.7}, we have
\[
\label{IV.3.3.9.1}
q(\operatorname{Ass}(\sh{F}\otimes_S\sh{G}))\subset\operatorname{Ass}(\sh{G})
\tag{3.3.9.1}
\]
(where $q:X\times_S Y\to Y$ is the canonical projection).
\end{corollary}

\begin{proof}
Indeed, if $(x,y)\in Z$, we have $q(\mathrm{I}_{x,y})=\{y\}\subset\operatorname{Ass}(\sh{G})$.
\end{proof}

\oldpage[IV-2]{46}

\subsection{Properties of the sheaves \texorpdfstring{$\sh{F}/t\sh{F}$}{F/tF}}
\label{subsection:IV.3.4}
\label{IV.3.4}

\begin{proposition}[3.4.1]
\phantomsection
\label{IV.3.4.1}
Let $X$ be a locally Noetherian prescheme, $t$ a section of $\sh{O}_X$ over $X$, and $Y$ the closed sub-prescheme of $X$ defined by the ideal $t\sh{O}_X$.
Let $\sh{F}$ be a coherent $\sh{O}_X$-module, let $S$ be the reduced closed sub-prescheme of $X$ having $\operatorname{Supp}(\sh{F})$ as underlying space, and let $(S_i)$ be the family of reduced closed sub-preschemes of $X$ whose underlying spaces are the irreducible components of $S$; denote by $s_i$ the generic point of $S_i$.
Finally, let $Z$ be an irreducible component of $\operatorname{Supp}(\sh{F}/t\sh{F})=S\cap Y$, $z$ its generic point.
\begin{enumerate}
  \item[\emph{(i)}] For every $i$ such that $Z\subset S_i$, $Z$ is an irreducible component of $S_i\cap Y$.
  \item[\emph{(ii)}] If $Z$ is not equal to any of the $S_i$, then
\[
\label{IV.3.4.1.1}
\operatorname{length}((\sh{F}/t\sh{F})_z)\geqslant\sum_i\operatorname{length}(\sh{F}_{s_i})
\tag{3.4.1.1}
\]
where the sum on the right-hand side is extended to all $i$ such that $Z\subset S_i$.
  \item[\emph{(iii)}] Suppose that $Z$ is not equal to any of the $S_i$.
  For the two sides of \sref{IV.3.4.1.1} to be equal, it is necessary and sufficient that the following two conditions be satisfied:
  \begin{enumerate}
    \item[\emph{$\alpha$)}] $t_z$ is $\sh{F}_z$-regular \sref[0]{0.15.1.4}.
    \item[\emph{$\beta$)}] For every $i$ such that $Z\subset S_i$, the canonical image of the germ $t_z$ in $\sh{O}_{S_i,z}$ generates the maximal ideal of this ring (which implies that $\sh{O}_{S_i,z}$ is a discrete valuation ring and that the image of $t_z$ is a uniformizer of this ring).
  \end{enumerate}
\end{enumerate}
\end{proposition}

\begin{proof}
\emph{(i)} If $j:Y\to X$ is the canonical injection, then
\[
\sh{F}/t\sh{F}
  =\sh{F}\otimes_{\sh{O}_X}\sh{O}_Y
  =j^*\sh{F},
\]
and hence
\[
\operatorname{Supp}(\sh{F}/t\sh{F})
  =j^{-1}(S)=S\cap Y
\]
by \sref[I]{I.9.1.13}, which proves the assertion.

\emph{(ii) and (iii)} The points $s_i$ such that $Z\subset S_i$ are precisely those belonging to $\operatorname{Spec}(\sh{O}_z)$.
Thus, to prove \emph{(ii)} and \emph{(iii)}, we may replace $X$ by $\operatorname{Spec}(\sh{O}_z)$.
If $M=\sh{F}_z$, we may therefore suppose that $\sh{F}=\widetilde{M}$, and then $\sh{F}_{s_i}=M_{\mathfrak p_i}$, where the $\mathfrak p_i$ are the minimal prime ideals of $\sh{O}_z$.
Moreover, since $M$ is an $\sh{O}_z$-module of finite type, we have $S=S'_{\mathrm{red}}$, where
\[
S'=\operatorname{Spec}(\sh{O}_z/\mathfrak a)
\]
and $\mathfrak a$ is the annihilator of $M$ in $\sh{O}_z$ \sref[0]{0.1.7.4}.
The two sides of \sref{IV.3.4.1.1} have the same values whether $M$ is regarded as an $\sh{O}_z$-module or as an $(\sh{O}_z/\mathfrak a)$-module.
We may therefore finally replace $X$ by $S'=\operatorname{Spec}(A)$, where $A$ is a Noetherian local ring and $M$ is a faithful $A$-module.
Since $Z$ is closed in $S$, the hypothesis that $Z\neq S_i$ for every $i$ means that $s_i\notin Z$, and hence that $\dim(A)>0$.
Finally, to say that $z$ is the generic point of $Z$, an irreducible component of $S\cap V(t)$, means that $A/tA$ has dimension $0$ (in other words, that it is a local Artinian ring).
We are thus reduced to the following lemma.
\end{proof}

\begin{lemma}[3.4.1.2]
\phantomsection
\label{IV.3.4.1.2}
Let $A$ be a Noetherian local ring of dimension $>0$, let $\mathfrak p_i$ be the minimal prime ideals of $A$, let $\mathfrak m$ be its maximal ideal, and let $t\in\mathfrak m$ be such that $A/tA$ is Artinian.
Then, for every $A$-module of finite type $M$,
\[
\label{IV.3.4.1.3}
\operatorname{length}(M/tM)
  \geqslant
  \sum_i\operatorname{length}(M_{\mathfrak p_i}).
\tag{3.4.1.3}
\]
Moreover, equality holds if and only if the following two conditions are satisfied:
\begin{enumerate}
  \item[\emph{$\alpha$)}] $t$ is $M$-regular;
  \item[\emph{$\beta$)}] for every $i$ such that $M_{\mathfrak p_i}\neq 0$, the image of $t$ in $A/\mathfrak p_i$ generates the maximal ideal of that ring (which implies that $A/\mathfrak p_i$ is a discrete valuation ring).
\end{enumerate}
\end{lemma}

\oldpage[IV-2]{47}

\begin{proof}
Since $A$ does not have dimension $0$ and $A/tA$ is Artinian, necessarily $\dim(A)=1$ \sref[0]{0.16.3.4}, and $t\notin\mathfrak p_i$ for every $i$.
The principal ideal $(t)$ is therefore an ideal of definition of $A$, and hence contains a power of the maximal ideal $\mathfrak m$.
Let $N$ be the submodule of $M$ consisting of the elements annihilated by a power of $t$ (or, equivalently in view of what we have just proved, by a power of $\mathfrak m$), and put $P=M/N$.
Then $t$ is $P$-regular: if $tx\in N$ for some $x\in M$, then $t^k(tx)=0$ for some integer $k$, and hence $x\in N$.
We shall use the following elementary lemma.

\begin{lemma}[3.4.1.4]
\phantomsection
\label{IV.3.4.1.4}
Let $A$ be a ring and
\[
0\longrightarrow M'\longrightarrow M\longrightarrow M''\longrightarrow 0
\]
an exact sequence of $A$-modules.
If $t\in A$ is $M''$-regular, then the sequence
\[
0\longrightarrow M'/tM'
 \longrightarrow M/tM
 \longrightarrow M''/tM''
 \longrightarrow 0
\]
is exact.
\end{lemma}

\begin{proof}
Since $M/tM=M\otimes_A(A/tA)$, it suffices to prove exactness at $M'/tM'$.
Suppose that the image $x\in M$ of an element $x'\in M'$ is of the form $x=ty$, with $y\in M$.
For the images $x'',y''$ of $x,y$ in $M''$, we have $x''=ty''$.
But $x''=0$, so the hypothesis gives $y''=0$.
Thus $y$ is the image of an element $y'\in M'$, and the relation $x=ty$ gives $x'=ty'$ because $M'\to M$ is injective.
\end{proof}

It follows from this lemma that
\[
\label{IV.3.4.1.5}
\operatorname{length}(M/tM)
  =
  \operatorname{length}(N/tN)
  +
  \operatorname{length}(P/tP).
\tag{3.4.1.5}
\]
On the other hand, $N_{\mathfrak p_i}=0$ for every $i$, since $t\notin\mathfrak p_i$, and hence $M_{\mathfrak p_i}=P_{\mathfrak p_i}$.
Thus, to prove \sref{IV.3.4.1.3}, it suffices to prove it with $P$ in place of $M$.
Furthermore, if the two sides of \sref{IV.3.4.1.3} are equal, then the same inequality for $P$, together with \sref{IV.3.4.1.5}, gives $\operatorname{length}(N/tN)=0$.
Consequently $N/tN=0$, and Nakayama's lemma gives $N=0$, since $N$ is of finite type.
But $N=0$ means that $t$ is $M$-regular.
We may therefore reduce to the case $M=P$, that is, suppose from the outset that $t$ is $M$-regular.
This implies that $\mathfrak m\notin\operatorname{Ass}(M)$, since $t$ cannot annihilate a nonzero element of $M$.
Because $\dim(A)=1$, it follows that
\[
\operatorname{Ass}(M)\subset\{\mathfrak p_i\}_i.
\]

We now argue by induction on
\[
n=\sum_i\operatorname{length}(M_{\mathfrak p_i}).
\]
If $n=0$, then $M_{\mathfrak p_i}=0$ for every $i$, and hence $M=0$, since none of the $\mathfrak p_i$ belongs to $\operatorname{Ass}(M)$.
Both sides of \sref{IV.3.4.1.3} are then zero, and condition \emph{$\beta$)} is vacuous.
If $n>0$, the argument at the beginning of the proof of \sref{IV.3.4.1} allows us, moreover, to suppose that $M$ is faithful.
It follows that $M_{\mathfrak p_i}\neq 0$ for every $i$ (Bourbaki, \emph{Alg.\ comm.}, chap.~II, \textsection~2, no~2, cor.~2 of prop.~4), and consequently
\[
\operatorname{Ass}(M)=\{\mathfrak p_i\}_i.
\]

\oldpage[IV-2]{48}

Suppose first that $n=1$.
There is then only one minimal prime ideal $\mathfrak p$ of $A$, and the assertion that $M_{\mathfrak p}$ has length $1$ means that it is isomorphic, as an $A_{\mathfrak p}$-module, to the residue field
\[
k=A_{\mathfrak p}/\mathfrak pA_{\mathfrak p}.
\]
Thus $M_{\mathfrak p}$ is annihilated by $\mathfrak pA_{\mathfrak p}$, and consequently $\mathfrak p$ is the annihilator of $M$ (Bourbaki, \emph{Alg.\ comm.}, chap.~II, \textsection~2, no~4, formula~(9)).
Since $M$ is faithful, $\mathfrak p=0$, and hence $A$ is integral.
Because $M\neq 0$, Nakayama's lemma gives $M/tM\neq 0$, so that $\operatorname{length}(M/tM)\geqslant 1$; this proves \sref{IV.3.4.1.3} in this case.
Moreover, if $\operatorname{length}(M/tM)=1$, then $M$ is cyclic (Bourbaki, \emph{Alg.\ comm.}, chap.~II, \textsection~3, no~2, cor.~2 of prop.~4), and hence isomorphic to a quotient $A/\mathfrak b$.
Since $M$ is faithful, necessarily $\mathfrak b=0$, and $M$ is isomorphic to $A$.
As $\operatorname{length}(A/tA)=1$, the ideal $tA$ is the maximal ideal $\mathfrak m$.
Since $A$ is a Noetherian integral local ring, it follows that $A$ is a discrete valuation ring (Bourbaki, \emph{Alg.\ comm.}, chap.~VI, \textsection~3, no~6, prop.~9), with uniformizer $t$.
Conversely, if $A$ is a discrete valuation ring, $t$ is its uniformizer, $\operatorname{length}(M_{\mathfrak p})=1$, and $t$ is $M$-regular, then $M$ is torsion-free and therefore isomorphic to a submodule of $A$ (as $M$ is of finite type), hence to $A$ itself.
Thus
\[
\operatorname{length}(M/tM)=\operatorname{length}(A/tA)=1.
\]

Suppose now that $n\geqslant 2$.
There exists an exact sequence
\[
0\longrightarrow M'\longrightarrow M\longrightarrow M''\longrightarrow 0
\]
with $M'\neq 0$, $M''\neq 0$, and
\[
\operatorname{Ass}(M)=\operatorname{Ass}(M')\cup\operatorname{Ass}(M'').
\]
Indeed, if $\operatorname{Ass}(M)$ has more than one element, this follows from Bourbaki, \emph{Alg.\ comm.}, chap.~IV, \textsection~1, no~1, prop.~4.
If, on the contrary, $\operatorname{Ass}(M)$ consists of a single prime ideal, that ideal is necessarily the unique minimal prime ideal $\mathfrak p$ of $A$.
The hypothesis then gives $\operatorname{length}(M_{\mathfrak p})\geqslant 2$, and it is enough to take for $M'$ the inverse image of a submodule of $M_{\mathfrak p}$ distinct from both $0$ and $M_{\mathfrak p}$.
Since $t$ is $M$-regular, it belongs to none of the prime ideals in $\operatorname{Ass}(M)$ (Bourbaki, \emph{Alg.\ comm.}, chap.~IV, \textsection~1, no~1, cor.~2 of prop.~2).
For the same reason, $t$ is both $M'$-regular and $M''$-regular.
By \sref{IV.3.4.1.4}, the sequence
\[
0\longrightarrow M'/tM'
 \longrightarrow M/tM
 \longrightarrow M''/tM''
 \longrightarrow 0
\]
is exact.
For every $i$, so is the sequence
\[
0\longrightarrow M'_{\mathfrak p_i}
 \longrightarrow M_{\mathfrak p_i}
 \longrightarrow M''_{\mathfrak p_i}
 \longrightarrow 0.
\]
Consequently
\[
\begin{aligned}
\operatorname{length}(M/tM)
  &=
  \operatorname{length}(M'/tM')
  +
  \operatorname{length}(M''/tM''),\\
\operatorname{length}(M_{\mathfrak p_i})
  &=
  \operatorname{length}(M'_{\mathfrak p_i})
  +
  \operatorname{length}(M''_{\mathfrak p_i}).
\end{aligned}
\]
The induction hypothesis now proves \sref{IV.3.4.1.3}.
Moreover, the two sides can be equal only if the analogous inequalities for $M'$ and $M''$ are both equalities.
By the induction hypothesis, this is equivalent to condition \emph{$\beta$)} for those $\mathfrak p_i$ such that $M'_{\mathfrak p_i}\neq 0$ or $M''_{\mathfrak p_i}\neq 0$; these are precisely the ideals for which $M_{\mathfrak p_i}\neq 0$.
\end{proof}

\oldpage[IV-2]{49}

\begin{corollary}[3.4.2]
\phantomsection
\label{IV.3.4.2}
Under the general hypotheses of \sref{IV.3.4.1}, suppose that $Z$ is not equal to any of the $S_i$ and that $\operatorname{length}((\sh{F}/t\sh{F})_z)=1$.
Then there is a unique $S_i$ containing $Z$, and for this value of $i$, $\operatorname{length}(\sh{F}_{s_i})=1$; moreover, $\sh{O}_{S_i,z}$ is a discrete valuation ring with uniformizer $t_z$, and $t_z$ is $\sh{F}_z$-regular.
\end{corollary}

\begin{proof}
This follows from \sref{IV.3.4.1}, the two sides of \sref{IV.3.4.1.1} being then equal.
\end{proof}

\begin{proposition}[3.4.3]
\phantomsection
\label{IV.3.4.3}
Let $X$ be a locally Noetherian prescheme, $t$ a section of $\sh{O}_X$ over $X$, and $Y$ the closed sub-prescheme of $X$ defined by the ideal $t\sh{O}_X$.
Let $\sh{F}$ be a coherent $\sh{O}_X$-module, let $T$ be an associated prime cycle of $\sh{F}$, let $T'$ be an irreducible component of $T\cap Y$, and let $x$ be the generic point of $T'$.
Suppose that $t_x$ is $\sh{F}_x$-regular; then $x\in\operatorname{Ass}(\sh{F}/t\sh{F})$.
\end{proposition}

\begin{proof}
As in the proof of \sref{IV.3.4.1}, we may reduce to the case $X=\operatorname{Spec}(\sh{O}_x)$.
Taking \sref{IV.3.1.2} into account, the proposition is then a consequence of \sref[0]{0.16.4.6.3}.
\end{proof}

\begin{proposition}[3.4.4]
\phantomsection
\label{IV.3.4.4}
Let $X$ be a locally Noetherian prescheme, $t$ a section of $\sh{O}_X$ over $X$, and $Y$ the closed sub-prescheme of $X$ defined by the ideal $t\sh{O}_X$.
Let $\sh{F}$ be a coherent $\sh{O}_X$-module, and let $(S_i)$ be the family of irreducible components of $\operatorname{Supp}(\sh{F})$.
Let $y$ be a point of $Y$ such that $t_y$ is $\sh{F}_y$-regular and no embedded associated prime cycle of $\sh{F}/t\sh{F}$ contains $y$.
Then the irreducible components of $\operatorname{Supp}(\sh{F}/t\sh{F})$ containing $y$ are precisely the irreducible components of the $S_i\cap Y$ containing $y$, and the associated prime cycles of $\sh{F}$ containing $y$ are non-embedded.
\end{proposition}

\begin{proof}
Let us first prove the last assertion.
Let $T\supset T_1$ be two associated prime cycles of $\sh{F}$ containing $y$.
If $x$ is the generic point of an irreducible component of $T\cap Y$ containing $y$, then $x$ is a generization of $y$, and hence belongs to every neighbourhood of $y$.
The hypothesis that $t_y$ is $\sh{F}_y$-regular therefore implies that $t_x$ is $\sh{F}_x$-regular \sref[0]{0.15.2.4}, so \sref{IV.3.4.3} gives $x\in\operatorname{Ass}(\sh{F}/t\sh{F})$.
Let $x_1$ be the generic point of an irreducible component of $T_1\cap Y$ containing $y$, and let $x$ be the generic point of an irreducible component of $T\cap Y$ containing $x_1$.
The preceding argument shows that $x_1$ and $x$ both belong to $\operatorname{Ass}(\sh{F}/t\sh{F})$.
Since $x_1\in\overline{\{x\}}$, the hypothesis implies $x_1=x$.

Again denote by $T$ and $T_1$ the integral closed sub-preschemes of $X$ having these spaces as their respective underlying spaces, and set
\[
A=\sh{O}_{T,x},
\qquad
A_1=\sh{O}_{T_1,x}.
\]
Then $A_1=A/\mathfrak p$ for a prime ideal $\mathfrak p$ of $A$.
By the definitions of $x$ and $x_1$, the rings $A/tA$ and $A_1/tA_1$ are Artinian.
On the other hand, we saw above that $t_x$ is $\sh{F}_x$-regular, so $x\notin\operatorname{Ass}(\sh{F})$, and consequently $A_1$ is not Artinian.
Thus $\dim A=\dim A_1=1$.
This implies $\mathfrak p=0$ and $A=A_1$ \sref[0]{0.16.1.2.2}.
Since $\operatorname{Spec}(A)$ and $\operatorname{Spec}(A_1)$ are respectively dense in $T$ and $T_1$, it follows that $T=T_1$.

The $S_i$ containing $y$ are therefore all the associated prime cycles of $\sh{F}$ containing $y$.
If $x$ is the generic point of an irreducible component of $S_i\cap Y$ containing $y$, then \sref[0]{0.15.2.4} again shows that $t_x$ is $\sh{F}_x$-regular, and \sref{IV.3.4.3} gives $x\in\operatorname{Ass}(\sh{F}/t\sh{F})$.
This proves the first assertion of \sref{IV.3.4.4}.
\end{proof}

\begin{proposition}[3.4.5]
\phantomsection
\label{IV.3.4.5}
Let $X$ be a locally Noetherian prescheme, $t$ a section of $\sh{O}_X$ over $X$, and $Y$ the closed sub-prescheme of $X$ defined by the ideal $t\sh{O}_X$.
Let $\sh{F}$ be a coherent $\sh{O}_X$-module and $y$ a point of $Y$.
\oldpage[IV-2]{50}
Suppose that $t_y$ is $\sh{F}_y$-regular and that $\sh{F}/t\sh{F}$ is integral at $y$ \sref{IV.3.2.4}.
Then $\sh{F}$ is integral at the point $y$.
\end{proposition}

\begin{proof}
Taking into account \sref{IV.3.4.4}, it suffices to prove that $y$ is contained in a \emph{unique} irreducible component of $\operatorname{Supp}(\sh{F})$, and that if $s$ is the generic point of this component, $\operatorname{length}(\sh{F}_s)=1$.
By hypothesis, $y$ belongs to only one irreducible component of $\operatorname{Supp}(\sh{F}/t\sh{F})$, and if $z$ is the generic point of this component, then $\operatorname{length}((\sh{F}/t\sh{F})_z)=1$.
The conclusion follows from \sref{IV.3.4.2}.
\end{proof}

\begin{proposition}[3.4.6]
\phantomsection
\label{IV.3.4.6}
The hypotheses being those of \sref{IV.3.4.1}, let $x$ be a point of $Y$.
Suppose that $Y$ does not contain any of the $S_i$ containing $x$, and that $\sh{F}/t\sh{F}$ is reduced at the point $x$ \sref{IV.3.2.2}.
Then $t_x$ is $\sh{F}_x$-regular and $\sh{F}$ is reduced at the point $x$.
Moreover, if $z_j$ is the generic point of an irreducible component of $\operatorname{Supp}(\sh{F}/t\sh{F})$ containing $x$, then $z_j$ is contained in a unique $S_i$, and $\sh{O}_{S_i,z_j}$ is a discrete valuation ring with uniformizer $t_{z_j}$.
\end{proposition}

\begin{proof}
The assertion that $t_x$ is $\sh{F}_x$-regular follows by applying the following lemma to the ring $\sh{O}_x$.

\begin{lemma}[3.4.6.1]
\phantomsection
\label{IV.3.4.6.1}
Let $A$ be a Noetherian ring, $M$ an $A$-module of finite type, $\mathfrak p_i$ the minimal elements of $\operatorname{Supp}(M)$, and $t\in A$.
Suppose that $t$ belongs to none of the $\mathfrak p_i$ and that $M/tM$ is a reduced $A$-module \sref{IV.3.2.2}.
Then $t$ is $M$-regular.
\end{lemma}

\begin{proof}
Every prime ideal $\mathfrak p\in\operatorname{Supp}(M)$ contains one of the $\mathfrak p_i$.
Since $t$ belongs to none of the $\mathfrak p_i$, multiplication by $t$ on $M_{\mathfrak p}$ is not nilpotent (Bourbaki, \emph{Alg.\ comm.}, chap.~IV, \textsection~1, no~4, cor.\ of prop.~9).
Let $N$ be the submodule of $M$ consisting of the elements annihilated by a power of $t$, and set $P=M/N$.
We shall prove that $N=0$.
Since $t$ is $P$-regular, \sref{IV.3.4.1.4} gives an exact sequence
\[
0\longrightarrow N/tN
 \longrightarrow M/tM
 \longrightarrow P/tP
 \longrightarrow 0.
\]
Since $N$ is of finite type, it is annihilated by a power of $t$, and it therefore suffices to prove that $N/tN=0$.
As $N/tN$ is a submodule of $M/tM$, it suffices to prove that $(N/tN)_{\mathfrak p}=0$ for every $\mathfrak p\in\operatorname{Ass}(M/tM)$, or equivalently that
\[
u_{\mathfrak p}:(M/tM)_{\mathfrak p}\longrightarrow(P/tP)_{\mathfrak p}
\]
is bijective for every $\mathfrak p\in\operatorname{Ass}(M/tM)$.
Now $(P/tP)_{\mathfrak p}\neq 0$.
Indeed,
\[
\mathfrak p\in\operatorname{Supp}(M/tM)
  =\operatorname{Supp}(M)\cap V(t)
\]
by \sref[0]{0.1.7.5}, so the image of $t$ in $A_{\mathfrak p}$ lies in the maximal ideal $\mathfrak pA_{\mathfrak p}$.
The equality $P_{\mathfrak p}=tP_{\mathfrak p}$ would therefore imply $P_{\mathfrak p}=0$ by Nakayama's lemma.
We would then have $M_{\mathfrak p}=N_{\mathfrak p}$, and multiplication by $t$ on $M_{\mathfrak p}$ would be nilpotent, contradicting the observation above because $\mathfrak p\in\operatorname{Supp}(M)$.
Finally, since $M/tM$ is reduced, $\operatorname{length}((M/tM)_{\mathfrak p})=1$.
As $(P/tP)_{\mathfrak p}\neq 0$, the map $u_{\mathfrak p}$ is necessarily bijective.
This proves the lemma.
\end{proof}

By hypothesis, no embedded associated prime cycle of $\sh{F}/t\sh{F}$ contains $x$.
It follows from \sref{IV.3.4.4} that no embedded associated prime cycle of $\sh{F}$ contains $x$.
On the other hand, applying \sref{IV.3.4.2} to an irreducible component of $\operatorname{Supp}(\sh{F}/t\sh{F})$ containing $x$, we find that $\operatorname{length}(\sh{F}_{s_i})=1$ for every $S_i$ containing $x$.
This completes the proof that $\sh{F}_x$ is reduced.
The final assertions also follow from \sref{IV.3.4.2}.
\end{proof}

\begin{corollary}[3.4.7]
\phantomsection
\label{IV.3.4.7}
Let $A$ be a Noetherian local ring, $\mathfrak m$ its maximal ideal, $M$ an $A$-module of finite type, and $(x_i)_{1\leqslant i\leqslant k}$ a family of elements of $\mathfrak m$ forming part of a system of parameters for $M$ \sref[0]{0.16.3.6}.
\oldpage[IV-2]{51}
If the $A$-module
\[
N=M\Big/\Big(\sum_{i=1}^k x_iM\Big)
\]
is integral \sref{IV.3.2.4}, then $M$ is integral and the sequence $(x_i)_{1\leqslant i\leqslant k}$ is $M$-regular.
\end{corollary}

\begin{proof}
Induction on $k$ immediately reduces us to the case $k=1$; write $x$ for $x_1$.
The hypothesis that $x$ forms part of a system of parameters for $M$ implies
\[
\dim(N)=\dim(M)-1
\]
by \sref[0]{0.16.3.7}.
Put $n=\dim(N)$.
There is therefore a minimal element $\mathfrak p$ of $\operatorname{Supp}(M)$ such that
\[
\dim(M/\mathfrak pM)=n+1
\]
by \sref[0]{0.16.3.4}; for every integer $j>0$ one also has
\[
\dim(M/\mathfrak p^jM)=n+1
\]
by \sref[0]{0.16.3.5}.
Moreover, $x$ forms part of a system of parameters for $M/\mathfrak p^jM$ \sref[0]{0.16.3.5}.
Set
\[
M'=M/\mathfrak p^jM,
\qquad
N'=M'/xM';
\]
then $\dim(N')=n$.

There is an evident surjective homomorphism $v:N\to N'$.
We claim that it is bijective.
Indeed, if $P=\operatorname{Ker}(v)$, then $\operatorname{Ass}(P)\subset\operatorname{Ass}(N)$.
Since $N$ is integral, if $P\neq 0$, then $\operatorname{Ass}(P)$ and $\operatorname{Ass}(N)$ would both consist of the unique point $\mathfrak q$ of
\[
\operatorname{Ass}(N)=\operatorname{Supp}(N).
\]
But $\dim(N')=\dim(N)$ implies $N'_{\mathfrak q}\neq 0$, and $\operatorname{length}(N_{\mathfrak q})=1$ implies $\operatorname{length}(N'_{\mathfrak q})=1$, since $N_{\mathfrak q}\to N'_{\mathfrak q}$ is surjective.
Thus $P_{\mathfrak q}=0$, a contradiction.
Hence $v$ is bijective.

It follows that $N'$, being isomorphic to $N$, is integral.
Moreover, the support of $N'$, which is the intersection of $\operatorname{Supp}(M')$ with $V(x)$, cannot contain $\operatorname{Supp}(M')$, and the latter is irreducible by construction.
Since $M'/xM'$ is integral, and hence reduced, \sref{IV.3.4.6} shows that $x$ is $M'$-regular.
Consequently the kernel of multiplication by $x$ on $M$ is contained in
\[
\mathfrak p^jM\subset\mathfrak m^jM
\]
for every integer $j$, and this kernel is therefore zero by \sref[0]{0.7.3.5}.
Thus $x$ is $M$-regular.
Applying \sref{IV.3.4.5}, we conclude that $M$ is integral.
\end{proof}

\begin{remark}[3.4.8]
\phantomsection
\label{IV.3.4.8}
The analogue of \sref{IV.3.4.7} obtained by replacing ``integral'' by ``reduced'' need not be true.
For example, let $C=K[X,Y,Z]$ be the polynomial ring over a field $K$, let
\[
B=C/\mathfrak p\mathfrak q,
\qquad
\mathfrak p=CZ,
\qquad
\mathfrak q=CX^2+CY,
\]
and let $A$ be the local ring of $B$ corresponding to the maximal ideal that is the image of $CX+CY+CZ$ in $B$.
If $x,z$ are the canonical images of $X,Z$ in $A$, then $xz\neq 0$ but $x^2z^2=0$.
On the other hand, $A/xA$ is isomorphic to $K[Y,Z]/(YZ)$, so $\dim(A/xA)=1$ whereas $\dim(A)=2$.
Thus $x$ belongs to a system of parameters of $A$ \sref[0]{0.16.3.4}, and $A/xA$ is reduced, but $A$ is not.
\end{remark}

\begin{proposition}[3.4.9]
\phantomsection
\label{IV.3.4.9}
Let $A$ be a Noetherian ring, $M$ an $A$-module of finite type, and $f\in A$ an $M$-regular element such that $M/fM$ has no embedded associated prime ideals.
If $\mathfrak p_i$ $(1\leqslant i\leqslant m)$ are the prime ideals associated to $M/fM$, then, for every integer $n>0$, $f^nM$ is the intersection of the inverse images of the $f^nM_{\mathfrak p_i}$ under the canonical maps $M\to M_{\mathfrak p_i}$ $(1\leqslant i\leqslant m)$.
\end{proposition}

\begin{proof}
It suffices to prove that the $\mathfrak p_i$ are also the prime ideals associated to $M/f^nM$, for then the $\mathfrak p_i$-saturations of $f^nM$ are the components of the necessarily unique reduced primary decomposition of $f^nM$ in $M$.
Now, we have
\[
\operatorname{Ass}(f^{n-1}M/f^nM)
\subset
\operatorname{Ass}(M/f^nM)
\subset
\operatorname{Ass}(M/f^{n-1}M)
\cup
\operatorname{Ass}(f^{n-1}M/f^nM)
\]
by \sref{IV.3.1.7}.
Since $f$ is $M$-regular, $f^{n-1}M/f^nM$ is isomorphic to $M/fM$, and the result follows by induction on $n$.
\end{proof}
